\bookchapter{第二十八章：花對錢買到心儀機——手機選購終極指南}

\chapterintro{買手機這件事，很多人會提前看一週評測、比一個月價格，結果最後還是稀裡糊塗花五六千買了臺自己其實用不上的旗艦。這章不列引數列——那些你在官網都能看到。我想做的，是給你一套“選機思維”，以後不管市場上出了什麼新手機，你都能自己判斷它適不適合你。}

\booksection{零、三步選機法：需求預算機型}

很多人在網上問“三千塊買什麼手機好”，這個問題其實沒法回答——因為不知道你買手機主要用來幹什麼。買手機的正確思考順序應該是：

\begin{tcbinfo}{選機三板斧}
\begin{enumerate}
  \item \textbf{先想清楚你最在乎什麼}：拍照？打遊戲？續航？螢幕？系統流暢？還是只需要能打電話發微信就行？把需求排個序。
  \item \textbf{再定預算上限}：不是"大概多少錢"，是"最多能接受花多少錢"。定一個硬上限，選機的時候不會越看越貴。
  \item \textbf{最後才挑機型}：拿著你的需求排序和預算，去匹配機型。
\end{enumerate}

\end{tcbinfo}

\booksection{一、核心引數怎麼看——不用全懂，看懂這五項就夠了}

引數列動輒幾十項資料，其實普通使用者只需要關注這五個維度：

\begin{center}
\begin{tabularx}{\linewidth}{|X|X|X|}
\hline
\textbf{維度} & \textbf{怎麼看} & \textbf{一句話總結} \tabularnewline
\hline
處理器（SoC） & 決定了手機能流暢用幾年。2026年的標準：驍龍8 Elite Gen5/天璣9500以上=頂級；驍龍8s Gen4/天璣8400=中高階；驍龍7+ Gen4/天璣7350=夠用三年。低於驍龍6系的建議慎重。 & 買新不買舊，SoC決定了你手機能用幾年不卡 \tabularnewline
螢幕 & OLED比LCD色彩好、更省電；重新整理率120Hz比60Hz滑屏流暢很多；解析度至少1080P以上。高頻PWM調光（1920Hz以上）對眼睛更友好。 & OLED+120Hz+高頻調光=好螢幕的及格線 \tabularnewline
攝像頭 & \textbf{不看畫素數看感測器尺寸}。底大一級壓死人，感測器尺寸比畫素重要得多。超廣角實用；長焦（3x以上光學變焦）拍遠處的東西；微距鏡頭多半是湊數的。 & 大底主攝+實用長焦>多顆湊數鏡頭 \tabularnewline
電池和充電 & 4500mAh是基礎及格線，5000mAh以上一天一充很穩。快充看瓦數：60W以上算快，100W以上飛快。無線充電方便但慢。 & 5000mAh+60W快充=不焦慮組合 \tabularnewline
系統更新承諾 & 決定了手機能安全用多久。蘋果一般給5-6年系統更新；三星旗艦給4年系統+5年安全更新；谷歌Pixel給7年；國產中端機更新承諾普遍偏短（2-3年）。 & 想用久一點，選系統更新承諾長的品牌 \tabularnewline
\hline
\end{tabularx}
\end{center}

\booksection{二、按預算分檔推薦（以2026年中市場為基準）}

以下推薦基於2026年中國內地市場行情，價格為官方起售價。具體價格會有浮動，以你入手時的實際價格為準。

\booksubsection{1000元以下：能打電話，能刷微信，夠用就行}

\textbf{適合誰}：學生、備用機、對手機要求不高、主要用於打電話和微信。

\textbf{能買到什麼}：Redmi Note系列、榮耀X系列、realme數字系列的中低配。處理器大概驍龍6系或天璣6000系列，螢幕是LCD為主，拍照能拍清楚但對焦慢，系統流暢度一般。

\textbf{選購建議}：這個價位看兩點：電池夠不夠大（選5000mAh以上的），記憶體夠不夠（選6GB+128GB起步，別買4GB的會卡）。紅米在這個價位優勢最明顯。

\booksubsection{1000-2000元：價效比最高，能用三四年不卡的甜點區}

\textbf{適合誰}：學生、預算有限但想用得久、對拍照和遊戲沒有極致追求。

\textbf{代表性機型}：Redmi Turbo系列、iQOO Z系列、榮耀X系列高配、一加Ace系列入門。

\textbf{選購建議}：選搭載驍龍7+ Gen3或天璣7300以上處理器的機型，配8GB+256GB記憶體。這個配置日常使用+偶爾打王者榮耀完全夠用，正常能用3-4年不卡。

\booksubsection{2000-3500元：中堅力量，效能和拍照都有明顯提升}

\textbf{適合誰}：想要不錯的拍照、打遊戲不卡、希望用得久的普通使用者。這是最多人買的價格區間。

\textbf{代表性機型}：小米數字系列（如小米17標準版降價後）、vivo S系列高配、OPPO Reno系列、榮耀數字系列、Redmi K系列。

\textbf{選購建議}：這個價位最大的特點是“水桶機”——沒有明顯短板。比起2000元以下的機子，主要提升在：相機明顯變好（有光學防抖了）、螢幕上了高素質OLED、快充更快（80W以上常見）、做工質感也更精緻。選12GB+256GB或12GB+512GB，用四五年沒問題。

\booksubsection{3500-5000元：準旗艦，除了極客玩家其他人完全夠用}

\textbf{適合誰}：對拍照有追求的使用者、想體驗頂級質感的使用者。

\textbf{代表性機型}：小米數字系列標準版、vivo X系列、OPPO Find X系列標準版、iPhone老款旗艦（如iPhone 16）。

\textbf{選購建議}：這個價位已經能買到去年的旗艦機型。我的建議是：\textbf{買去年的旗艦，不買今年的中端}。去年旗艦的用料、做工、攝像頭硬體通常比今年的中端好一個檔次。前提是你能接受它不是“最新款”。

\booksubsection{5000元以上：旗艦區，不差錢或對某項體驗有極致追求}

\textbf{適合誰}：攝影發燒友、重度遊戲玩家、追求頂級螢幕和工藝的使用者、用手機拍影片做內容的創作者。

\textbf{代表性機型}：iPhone 17 Pro/Pro Max、三星S26 Ultra、華為Pura 80 Ultra、小米17 Ultra、vivo X300 Pro。

\textbf{選購建議}：這個價位各家的差距越來越小，其實主要看你喜歡什麼系統。喜歡拍照選vivo華為，喜歡拍影片選iPhone，喜歡螢幕和SPen選三星，喜歡均衡選小米。說實話到了這個價位，\textbf{你可以任性，選你看著順眼的那臺就行}。

\booksection{三、按場景選機——你的需求是什麼？}

以下是幾個最常見的“使用畫像”，對號入座就行：

\begin{center}
\begin{tabularx}{\linewidth}{|X|X|X|}
\hline
\textbf{使用場景} & \textbf{重點看什麼} & \textbf{推薦方向} \tabularnewline
\hline
拍照達人 & 大底主攝（1英寸最佳）、光學防抖、長焦（3x以上光學變焦）、影像晶片 & vivo X系列、華為Pura/Mate系列、小米Ultra系列 \tabularnewline
遊戲玩家 & 頂級處理器、高觸控取樣率、散熱系統、大電池+快充 & iQOO數字系列、ROG Phone、紅魔、Redmi K系列 \tabularnewline
續航優先 & 5500mAh以上電池、能效好的處理器、LTPO螢幕 & 榮耀Magic系列、一加、Redmi Note大電池版 \tabularnewline
長輩使用 & 簡易模式、大字型、大音量、大電池、售後網點多 & 華為暢享、榮耀Play、OPPO A系列、Redmi Note \tabularnewline
學生用機 & 價效比高、耐用、電池大、有家長控制功能 & Redmi Note/Turbo、realme GT、iQOO Neo \tabularnewline
蘋果生態使用者 & 已有Mac/iPad/AirPods，買iPhone是唯一解 & iPhone 17標準版（價效比）或Pro版（滿血體驗） \tabularnewline
出國旅行 & eSIM支援、全球頻段覆蓋完整、Google服務 & iPhone、三星S系列、Pixel、一加 \tabularnewline
\hline
\end{tabularx}
\end{center}

\booksection{四、五個容易踩的坑——買之前對照看一眼}

\begin{enumerate}
  \item \textbf{坑一：四攝五攝不等於拍照好}——廠商喜歡在引數上堆鏡頭數量，裡面經常有兩顆是200萬畫素的"湊數鏡頭"（微距+景深），實際完全沒用。看手機拍照能力，\textbf{只看主攝+超廣角+長焦}這三顆就夠。其他多出來的可以當它不存在。
  \item \textbf{坑二：不要被1億畫素唬住}——畫素高不代表畫質好。感測器尺寸、光圈、演演算法調教比畫素重要得多。很多旗艦機用5000萬畫素的大底感測器，拍出來的照片比某些1億畫素的千元機好太多。
  \item \textbf{坑三：別買運營商定製機}——"存電電話費送手機"看起來很劃算，但運營商定製機系統更新慢、可能有刪除不了的預裝App、而且通常是運營商鎖定的（只能用這家運營商的卡）。看似省了錢，實則"鎖定"了你。
  \item \textbf{坑四："曲面屏很高階"——是營銷，不是真理}——曲面屏曾經是旗艦機的標緻，但現在越來越多人發現它的問題：貼膜難、邊緣誤觸、碎了維修貴。2026年的趨勢是迴歸"微曲"甚至"純直屏"。不要因為"曲面的就是高階的"而選曲面屏。
  \item \textbf{坑五：不是所有"AI手機"都真的有用}——2026年所有廠商都在喊"AI手機"，但很多隻是把已有的語音助手換了個名字。真正有用的AI功能是：實時翻譯、AI消除（拍照把人去掉）、AI修圖、端側大模型推理、本地化的智慧摘要。其他的花哨功能基本都是噱頭。
\end{enumerate}

\booksection{五、什麼時候買最劃算？}

\begin{itemize}
  \item \textbf{618（5月底到6月中旬）}：一年中手機降價幅度最大的時段。很多手機廠商為了沖銷量，旗艦機能降價800-1500元。是上半年買手機的最佳視窗。
  \item \textbf{雙11（10月底到11月中旬）}：下半年的大促，力度和618差不多。這個時間段很多品牌新一代旗艦剛釋出，上一代旗艦降價幅度最大。
  \item \textbf{新機發布後一個月}：剛釋出時原價，一個月後渠道價（第三方店鋪）一般會便宜200-500。
  \item \textbf{春節前兩週到元宵節}：線下實體店清庫存的"年貨節"，可以和店員談價。
  \item \textbf{不推薦的時間}：新品釋出會後當天衝首發——除非你真的特別想要當第一批使用者，否則首發價是最貴的。
\end{itemize}

\booksection{六、渠道怎麼選——在哪買靠譜？}

\begin{itemize}
  \item \textbf{品牌官方商城/天貓旗艦店/京東自營}：最保險，價格沒優勢但不用擔心翻新和假貨，售後最好。
  \item \textbf{拼多多"百億補貼"}：價格最低，但需要選標註了"百億補貼"的店鋪。沒標註的不要碰。收到貨記住錄開箱影片，有問題能退貨。
  \item \textbf{線下實體店}：可以親手摸、試手感。注意：導購會推利潤高的機型（而不是最適合你的），建議進店之前先在網上鎖定2-3款目標機型，去了就說"我只想看這幾款"。
  \item \textbf{二手平臺（閒魚/轉轉）}：便宜但有風險。如果要買二手，記住三條：要求提供購買發票或訂單截圖、要求提供IMEI號查驗、\textbf{只走平臺交易，絕不私下轉帳}。
\end{itemize}

\begin{tcbpractice}{動手：用這套方法選一臺手機}
\begin{enumerate}
  \item 拿一張紙，寫下你的兩個核心需求（比如：續航第一、拍照第二）。
  \item 定一個預算上限（比如3500元）。
  \item 開啟京東或天貓，按價格篩選，圈出3-5款備選機型。
  \item 去B站搜"XX機型 一個月使用體驗"（不是第一天評測，是用了至少一個月之後的真實感受），看兩三個影片。
  \item 確認渠道靠譜後下單，收貨時錄開箱影片。
\end{enumerate}

\end{tcbpractice}

\begin{tcbsummary}{本章小結}
\begin{itemize}
  \item 選機三板斧：先想需求 再定預算 最後挑機型
  \item 核心引數只看五項：處理器、螢幕、攝像頭、電池充電、系統更新承諾
  \item 1500-3500元是價效比最高、最推薦的甜點區間
  \item 買去年的旗艦比買今年的中端更劃算
  \item 五大坑：多攝騙局、高畫素陷阱、運營商定製、曲面屏迷信、AI噱頭
  \item 618和雙11是一年中最適合買手機的兩個時間視窗
\end{itemize}

\end{tcbsummary}
